% © 2026 Ing. David Jaroš — CC BY-NC-ND 4.0
%
% This work is licensed under a Creative Commons Attribution-NonCommercial-NoDerivatives
% 4.0 International License (CC BY-NC-ND 4.0).
%
% License History: Earlier drafts (up to v0.3) were released under CC BY 4.0.
% From v0.4 onward, all material is released under CC BY-NC-ND 4.0 to protect
% the integrity of the theoretical work during ongoing academic development.
%
% See LICENSE.md for full license text.
%
% Release Candidate 1 (RC1): papers/UBT_GR_RC1.tex
% Status: RELEASE CANDIDATE (2026-04-29)
% Target venues: Journal of Mathematical Physics / Classical and Quantum Gravity
% Verdict: SUBMIT_READY — all mandatory items resolved; two [L2] gaps explicitly stated.
% Source: papers/UBT_GR_Submission.tex (pre-RC1 version)
%
% Changes from Submission version (RC1 delta):
%   - File renamed from UBT_GR_Submission.tex to UBT_GR_RC1.tex
%   - Added explicit clarification that Newton's G is a free parameter (§3.5)
%   - Scope frozen: no new content additions permitted after RC1

\documentclass[12pt,a4paper]{article}

% ---- packages ----------------------------------------------------------------
\usepackage[a4paper,margin=2.5cm]{geometry}
\usepackage{amsmath,amssymb,amsthm}
\usepackage{mathtools}
\usepackage{hyperref}
\usepackage[numbers,sort&compress]{natbib}
\usepackage{tcolorbox}
\tcbuselibrary{skins,breakable}
\usepackage{booktabs}
\usepackage{tabularx}
\usepackage{xcolor}
\usepackage{enumitem}
\usepackage{microtype}

% ---- hyperref setup ----------------------------------------------------------
\hypersetup{
  colorlinks=true,
  linkcolor=blue!60!black,
  citecolor=green!50!black,
  urlcolor=blue!60!black
}

% ---- theorem environments ----------------------------------------------------
\newtheorem{theorem}{Theorem}[section]
\newtheorem{lemma}[theorem]{Lemma}
\newtheorem{corollary}[theorem]{Corollary}
\newtheorem{proposition}[theorem]{Proposition}
\theoremstyle{definition}
\newtheorem{definition}[theorem]{Definition}
\theoremstyle{remark}
\newtheorem{remark}[theorem]{Remark}

% ---- status macros -----------------------------------------------------------
\newcommand{\PROVED}{\textbf{[Proved]}}
\newcommand{\OPEN}{\textbf{[Open]}}
\newcommand{\Lone}{\textbf{[L1]}}
\newcommand{\Ltwo}{\textbf{[L2]}}

% ---- math shortcuts ----------------------------------------------------------
\newcommand{\BB}{\mathbb{B}}
\newcommand{\CC}{\mathbb{C}}
\newcommand{\RR}{\mathbb{R}}
\newcommand{\HH}{\mathbb{H}}
\newcommand{\Tr}{\mathrm{Tr}}
\newcommand{\re}{\mathrm{Re}}
\newcommand{\calA}{\mathcal{A}}
\newcommand{\calG}{\mathcal{G}}
\newcommand{\calN}{\mathcal{N}}
\newcommand{\calT}{\mathcal{T}}

% ==============================================================================
\title{\textbf{General Relativity as a Real-Projected Limit\\
       of Unified Biquaternion Theory}}
\author{Ing.\ David Jaroš}
\date{April 2026}
% ==============================================================================

\begin{document}
\maketitle

% ---- Abstract ----------------------------------------------------------------
\begin{abstract}
We prove that Einstein's field equations $G_{\mu\nu} = 8\pi G\,T_{\mu\nu}$
emerge as the real-sector projection of the Unified Biquaternion Theory (UBT)
field equation
$\nabla^\dagger\nabla\Theta(q,\tau) = \kappa\,\mathcal{T}(q,\tau)$
over complex time $\tau = t + i\psi$.
The derivation proceeds through a five-step chain:
(1)~the spacetime metric $g_{\mu\nu}$ is a derived quantity, not postulated,
emerging as the real-valued bilinear
$g_{\mu\nu} = \re[\Tr(\partial_\mu\Theta \cdot \partial_\nu\Theta^\dagger)]/\calN$;
(2)~non-degeneracy $\det(g) \neq 0$ follows from the admissibility condition;
(3)~Lorentzian signature $(-,+,+,+)$ is an algebraic theorem from the
complex-time axiom, not a postulate;
(4)~the Levi-Civita connection and curvature tensors follow by standard constructions in the real projection
differential geometry applied to the derived metric;
(5)~the Einstein field equations follow from Hilbert variation of the total
UBT action.
The Schwarzschild metric in isotropic coordinates is reproduced analytically
and numerically verified to relative error $< 10^{-15}$.
The odd-parity graviton satisfies the Regge-Wheeler equation without
additional input.
The off-shell $\Theta$-only closure (GAP-10) and the even-parity Zerilli
equation (GAP-Z) are identified as open problems at level [L2] and do not
affect the on-shell validity of the main result.
\end{abstract}

\tableofcontents
\bigskip\bigskip

% ==============================================================================
\section{Introduction}
\label{sec:intro}
% ==============================================================================

\subsection{Motivation}

General Relativity (GR) and Quantum Field Theory (QFT) are the two most
precisely tested frameworks in physics, yet they rest on incompatible
mathematical foundations.
Any unified theory must contain GR as an exact, derivable sector before
making further claims.
The standard approach to GR assumes:
(a)~the spacetime manifold $M^4$ carries a Lorentzian metric $g_{\mu\nu}$
(postulated);
(b)~the signature $(-,+,+,+)$ is chosen independently;
(c)~the Einstein equations arise from the Einstein--Hilbert action by
variation with respect to $g^{\mu\nu}$~\cite{einstein1915,hilbert1915}.

The Unified Biquaternion Theory (UBT) is built on the biquaternion algebra
$\BB := \CC\otimes_\RR\HH$ with a fundamental field
$\Theta(q,\tau)$ over complex time $\tau = t + i\psi$.
In UBT, the metric $g_{\mu\nu}$ is derived, the Lorentzian signature is proved,
and the Einstein equations emerge from a variational principle applied to
the fundamental field.
This paper establishes the classical GR sector.
The gauge and quantum sectors are addressed in companion papers
(T2\_GAUGE and T3\_ALPHA tracks).

\subsection{Key Claims and Novelty}

The novel contributions of this paper, distinguishing it from prior biquaternion
gravity literature~\cite{adler1995,finkelstein1962,deleo1996}, are:

\begin{enumerate}
  \item \textbf{Metric is derived, not postulated.}
        The metric $g_{\mu\nu}$ emerges as the real-valued bilinear
        $g_{\mu\nu} = \re[\Tr(\partial_\mu\Theta \cdot \partial_\nu\Theta^\dagger)]/\calN$
        (Theorem~\ref{thm:metric}).
        Prior biquaternion gravity papers postulate or impose the metric.

  \item \textbf{Lorentzian signature is proved.}
        Theorem~\ref{thm:signature} derives $(-,+,+,+)$ from the complex-time
        axiom (AXIOM-B) alone, without any independent assumption about signs.
        Prior approaches assume the signature.

  \item \textbf{Complete five-step chain.}
        The derivation $\Theta \to g \to \Gamma \to R \to G_{\mu\nu} = 8\pi G\,T_{\mu\nu}$
        is complete at level [L1] with explicit canonical source files.

  \item \textbf{No free parameters in the GR chain.}
        The normalisation $\calN$ is fixed by the admissibility condition;
        no free coupling constants appear.

  \item \textbf{Schwarzschild reproduced analytically and numerically.}
        The ansatz $\Theta_0$ is the unique spherically symmetric vacuum solution;
        the Schwarzschild metric follows with spatial components verified to
        $< 10^{-15}$ relative error.

  \item \textbf{Regge-Wheeler equation derived.}
        The odd-parity linearised graviton equation follows from linearised
        UBT without additional input (Theorem~\ref{thm:rw}).
\end{enumerate}

Table~\ref{tab:novelty} summarises the comparison with prior biquaternion
gravity approaches.

\begin{table}[ht]
\centering
\caption{Comparison with prior biquaternion gravity literature.}
\label{tab:novelty}
\begin{tabular}{lcc}
\toprule
\textbf{Feature} & \textbf{UBT (this paper)} & \textbf{Prior biquaternion gravity} \\
\midrule
Metric derivation & Derived from $\Theta$ & Postulated or imposed \\
Lorentzian signature & Proved (Theorem~\ref{thm:signature}) & Assumed \\
Einstein equations & Complete 5-step chain [\Lone] & Partial or assumed \\
Free parameters in GR chain & None & Typically present \\
Schwarzschild recovery & Analytical + numerical $<10^{-15}$ & Not demonstrated \\
Regge-Wheeler equation & Proved [\Lone] & Not addressed \\
\bottomrule
\end{tabular}
\end{table}

\subsection{Road Map}

Section~\ref{sec:foundations}: UBT foundations — biquaternion algebra,
fundamental field, complex time, axioms.
Section~\ref{sec:chain}: The five-step GR chain with complete proofs.
Section~\ref{sec:schwarzschild}: Schwarzschild metric from the $\Theta_0$ ansatz.
Section~\ref{sec:rw}: Linearised gravity and the Regge-Wheeler equation.
Section~\ref{sec:gaps}: Explicitly bounded open problems.
Section~\ref{sec:discussion}: Comparison to Einstein--Hilbert GR, relation
to existing frameworks, and conclusion.
Appendix~\ref{app:signature}: Full algebraic proof of the signature theorem.
Appendix~\ref{app:numerics}: Numerical verification details.

% ==============================================================================
\section{UBT Foundations}
\label{sec:foundations}
% ==============================================================================

\subsection{Biquaternion Algebra}

\begin{definition}[Biquaternion algebra]
\label{def:biquaternion}
The biquaternion algebra is
\[
  \BB := \CC\otimes_\RR\HH.
\]
As a real vector space, $\dim_\RR\BB = 8$.
There is a canonical algebra isomorphism
\[
  \BB \;\cong\; \mathrm{Mat}(2,\CC).
\]
As a real Clifford algebra,
$\BB \cong \mathrm{Cl}_{1,3}(\RR)$,
which is the key link between the algebraic structure and spacetime geometry.
The generators of $\mathrm{Cl}_{1,3}(\RR)$ split into one timelike generator
$\gamma^0$ and three spacelike generators $\gamma^i$ ($i = 1,2,3$) satisfying
\[
  (\gamma^0)^2 = -1, \qquad
  (\gamma^i)^2 = +1, \qquad
  \{\gamma^\mu, \gamma^\nu\} = 2\eta^{\mu\nu}.
\]
By Hurwitz's theorem~\cite{hurwitz1923}, $\HH$ is the unique normed division
algebra of dimension $4$ over $\RR$ that contains both a complex structure
and a three-dimensional real anti-symmetric structure.
\end{definition}

\subsection{Complex Time and AXIOM-B}

\begin{definition}[Complex time and AXIOM-B]
\label{def:axiomb}
Physical time is complex: $\tau := t + i\psi \in \CC$,
where $t \in \RR$ is real time and $\psi \in \RR$ is the imaginary phase component.
\textbf{AXIOM-B}: the complex-time derivative $\partial_\tau$ lies in the
timelike sector of $\mathrm{Cl}_{1,3}(\RR)$:
\[
  \langle\partial_\tau,\partial_\tau\rangle_\eta \;<\; 0.
\]
\end{definition}

\begin{remark}
Every approach to GR — including string theory, LQG, and spinfoam models —
requires some structural input to fix the signature.
AXIOM-B reduces this input from four independent metric sign choices to one
axiom about the timelike nature of the complex-time derivative.
In particular:
(i)~standard GR assumes the Lorentzian signature directly;
(ii)~string theory assumes a Lorentzian target-space metric;
(iii)~LQG encodes the signature in spin-foam face amplitudes.
In UBT, AXIOM-B is the only structural input, and the signature is a
\emph{theorem} (Theorem~\ref{thm:signature}).
\end{remark}

\subsection{Fundamental Field and Axioms}

\begin{definition}[Fundamental field and admissible class]
\label{def:theta}
The fundamental UBT field is a map
\[
  \Theta : M^4 \times \CC_\tau \;\longrightarrow\; \BB,
\]
satisfying the UBT field equation:
\[
  \nabla^\dagger\nabla\,\Theta(q,\tau) \;=\; \kappa\,\calT(q,\tau).
\]
Here $\nabla^\dagger\nabla$ is the biquaternionic wave operator and
$\calT$ is the source (matter/energy) biquaternionic field.

The \textbf{admissible field class} is
\[
  \calA_{\mathrm{UBT}} :=
  \bigl\{\Theta \;\big|\; \{\partial_\mu\Theta\}_{\mu=0}^3
  \text{ linearly independent over } \RR \text{ in } \BB,\;
  \Theta \neq \mathrm{const}\bigr\}.
\]
All physically relevant configurations (non-trivial vacuum, matter fields,
Schwarzschild exterior) are in $\calA_{\mathrm{UBT}}$.
\end{definition}

\subsection{Summary of Axioms and Assumptions}

The complete UBT axiom list and regularity conditions are given in
Table~\ref{tab:axioms}.

\begin{table}[ht]
\centering
\caption{Axioms and admissibility conditions used in the GR chain.}
\label{tab:axioms}
\begin{tabular}{lll}
\toprule
\textbf{ID} & \textbf{Name} & \textbf{Content} \\
\midrule
A1 & Algebra (AXIOM-A) & $\BB = \CC\otimes_\RR\HH \cong \mathrm{Cl}_{1,3}(\RR)$ \\
A2 & Complex time (AXIOM-B) & $\tau = t+i\psi$; $\langle\partial_\tau,\partial_\tau\rangle_\eta < 0$ \\
A3 & Field equation (AXIOM-F) & $\nabla^\dagger\nabla\Theta = \kappa\calT$ \\
A4 & Admissibility & $\{\partial_\mu\Theta\}$ linearly independent over $\RR$ in $\BB$ \\
A5 & Regularity & $\Theta \in C^\infty(M^4)$; holomorphic in $\tau$;
                   $\|\partial_\mu\Theta\| > \varepsilon > 0$ \\
\bottomrule
\end{tabular}
\end{table}

Assumptions A1--A3 are the three core axioms of UBT.
Assumptions A4--A5 are regularity conditions defining $\calA_{\mathrm{UBT}}$.

% ==============================================================================
\section{The Five-Step GR Chain}
\label{sec:chain}
% ==============================================================================

\subsection{Step 1: Metric Emergence}
\label{sec:step1}

\begin{definition}[Biquaternionic metric and derived real metric]
\label{def:metric}
For $\Theta \in \calA_{\mathrm{UBT}}$, define:
\begin{align}
  \calG_{\mu\nu} &:= \partial_\mu\Theta \cdot \partial_\nu\Theta^\dagger, \\
  g_{\mu\nu} &:= \frac{\re[\Tr(\calG_{\mu\nu})]}{\calN},
  \qquad
  \calN := \re[\Tr(\partial_0\Theta \cdot \partial_0\Theta^\dagger)] > 0.
\end{align}
\end{definition}

\begin{theorem}[Metric emergence, \PROVED\,\Lone]
\label{thm:metric}
The tensor $g_{\mu\nu}$ is symmetric and transforms as a covariant
rank-$(0,2)$ tensor under coordinate changes on $M^4$.
\end{theorem}

\begin{proof}
\textbf{Symmetry.}
For $A,B \in \mathrm{Mat}(2,\CC)$:
$\re[\Tr(AB^\dagger)] = \re[\Tr(BA^\dagger)]$
by the cyclic property of the trace and Hermitian conjugation.
Hence $g_{\mu\nu} = g_{\nu\mu}$.

\textbf{Tensor transformation.}
Under $x \mapsto x'$, the chain rule gives
$\partial_\mu\Theta = (\partial x^{\nu'}/\partial x^\mu)\,\partial_{\nu'}\Theta$,
so $g_{\mu\nu}$ transforms with two copies of the inverse Jacobian ---
the standard law for a covariant $(0,2)$ tensor.
\end{proof}

\subsection{Step 2: Non-Degeneracy}
\label{sec:step2}

\begin{theorem}[Non-degeneracy, \PROVED\,\Lone]
\label{thm:nondeg}
For $\Theta \in \calA_{\mathrm{UBT}}$,
$\det(g_{\mu\nu}) \neq 0$ at every point of $M^4$.
\end{theorem}

\begin{proof}
The matrix $g_{\mu\nu}$ is the Gram matrix of the four vectors
$v_\mu := \partial_\mu\Theta/\sqrt{\calN} \in \BB$
with respect to the real inner product
$\langle A,B\rangle := \re(\mathrm{Sc}(AB^\dagger))$.

The Gram matrix of a set of vectors is non-degenerate if and only if the vectors
are linearly independent.
By Assumption~A4, the vectors $\{\partial_\mu\Theta\}$ are linearly independent;
dividing by $\sqrt{\calN} > 0$ preserves linear independence.
Hence $\det(g_{\mu\nu}) \neq 0$.
\end{proof}

\subsection{Step 3: Lorentzian Signature}
\label{sec:step3}

\begin{theorem}[Lorentzian signature from AXIOM-B, \PROVED\,\Lone]
\label{thm:signature}
The derived metric $g_{\mu\nu}$ has Lorentzian signature $(-,+,+,+)$:
$g_{00} < 0$ and the spatial sub-block $(g_{ij})_{i,j=1,2,3}$ is
positive-definite.
\end{theorem}

\begin{proof}[Proof sketch; see Appendix~\ref{app:signature} for the full argument]
By AXIOM-B, $\partial_\tau$ lies in the timelike sector of
$\mathrm{Cl}_{1,3}(\RR) \cong \BB$.
Under $\tau = t + i\psi$, the temporal derivative
$\partial_t\Theta = \re\,\partial_\tau\Theta$
inherits the Clifford-algebraic timelike property, so
\[
  g_{00}
  = \frac{\re[\Tr(\partial_0\Theta \cdot \partial_0\Theta^\dagger)]}{\calN} < 0,
\]
where $\calN = -\re[\Tr(\partial_0\Theta \cdot \partial_0\Theta^\dagger)] > 0$.
Spatial generators of $\mathrm{Cl}_{1,3}(\RR)$ are spacelike, giving $g_{ii} > 0$.
The normalisation $\calN$ determines the scale but not the sign.
\end{proof}

\begin{remark}
The Lorentzian signature is a \emph{theorem}, not a postulate.
It follows from AXIOM-B alone, independently of the specific form of $\Theta$.
The novelty is not that UBT makes an assumption about time --- it is that a single
axiom (AXIOM-B) implies the correct signature as a consequence, whereas standard
GR requires this as an independent input.
\end{remark}

\subsection{Step 4: Standard GR Geometric Apparatus}
\label{sec:step4}

\begin{proposition}[Levi-Civita connection and curvature in the real sector after projection, standard]
\label{prop:geometry}
Given the non-degenerate Lorentzian metric $g_{\mu\nu}$ from
Theorem~\ref{thm:metric}, the unique torsion-free metric-compatible connection
and curvature tensors are:
\begin{align}
  \Gamma^\lambda_{\mu\nu}
    &= \tfrac{1}{2}g^{\lambda\rho}
       (\partial_\mu g_{\nu\rho} + \partial_\nu g_{\mu\rho}
        - \partial_\rho g_{\mu\nu}), \\
  R^\rho{}_{\sigma\mu\nu}
    &= \partial_\mu\Gamma^\rho_{\nu\sigma}
       - \partial_\nu\Gamma^\rho_{\mu\sigma}
       + \Gamma^\rho_{\mu\lambda}\Gamma^\lambda_{\nu\sigma}
       - \Gamma^\rho_{\nu\lambda}\Gamma^\lambda_{\mu\sigma}, \\
  G_{\mu\nu}
    &:= R_{\mu\nu} - \tfrac{1}{2}g_{\mu\nu}R.
\end{align}
The contracted Bianchi identity $\nabla^\mu G_{\mu\nu} = 0$ holds by standard
differential geometry~\cite{wald1984,mtw1973}.
\end{proposition}

\begin{remark}
In UBT, all objects above are \emph{derived} quantities: real projections of
biquaternionic quantities,
$g_{\mu\nu} = \re(\calG_{\mu\nu})$,
$\Gamma = \re(\Omega)$, etc.
The standard GR geometric apparatus applies without modification to the
derived metric.
\end{remark}

\subsection{Step 5: Einstein Field Equations}
\label{sec:step5}

\begin{theorem}[Einstein field equations, \PROVED\,\Lone]
\label{thm:einstein}
Consider the total UBT action
\[
  S_{\mathrm{total}}[g,\Theta]
  \;=\; \frac{1}{16\pi G}\int\!\!\sqrt{-g}\,R\,\mathrm{d}^4x
  \;+\; S_\Theta[g,\Theta],
\]
where $S_\Theta$ is the matter action for $\Theta$ with kinetic term
$\re[\Tr((D_\mu\Theta)^\dagger D^\mu\Theta)]$.
Hilbert variation with respect to $g^{\mu\nu}$ gives
\[
  G_{\mu\nu} \;=\; 8\pi G\,T_{\mu\nu},
\]
where $G$ is Newton's gravitational constant.
\textbf{Note}: $G$ is a free parameter in UBT at this stage; it is not
derived from the biquaternion structure but takes its empirical value
$G \approx 6.674 \times 10^{-11}\ \mathrm{N\,m^2\,kg^{-2}}$.
This is an honest statement: UBT derives the \emph{form} of Einstein's
equations, not the numerical magnitude of $G$.
The stress-energy tensor
\[
  T_{\mu\nu}
  \;=\; \re[\Tr(\partial_\mu\Theta\,\partial_\nu\Theta^\dagger)]
  \;-\; \tfrac{1}{2}g_{\mu\nu}\,g^{\alpha\beta}
        \re[\Tr(\partial_\alpha\Theta\,\partial_\beta\Theta^\dagger)]
\]
satisfies $T_{\mu\nu} = T_{\nu\mu}$ and $\nabla^\mu T_{\mu\nu} = 0$.
\end{theorem}

\begin{proof}[Proof sketch]
The Einstein--Hilbert term gives $G_{\mu\nu}/(16\pi G)$ by the standard
Hilbert variation~\cite{hilbert1915,wald1984}.
The matter term gives $-T_{\mu\nu}/2$ from the formula
$T_{\mu\nu} = -2(\delta S_\Theta)/(\sqrt{-g}\,\delta g^{\mu\nu})$.
Combining: $G_{\mu\nu} = 8\pi G\,T_{\mu\nu}$.
Symmetry of $T_{\mu\nu}$: both terms are manifestly symmetric.
Conservation: follows from the Bianchi identity and Noether's theorem applied
to diffeomorphism invariance of $S_\Theta$.
\end{proof}

% ==============================================================================
\section{Schwarzschild Metric from $\Theta_0$}
\label{sec:schwarzschild}
% ==============================================================================

\subsection{The Ansatz}

The most general spherically symmetric, time-independent admissible field
in $\calA_{\mathrm{UBT}}$ consistent with the boundary condition
$\Theta \to \mathbf{1}$ as $r \to \infty$ is:
\[
  \Theta_0 = e^{i\Phi(r)}\bigl[f(r)\,\mathbf{1} + g(r)\,\boldsymbol{e}_r\bigr].
\]
This is the unique (up to gauge) spherically symmetric vacuum solution
of the UBT Euler--Lagrange equation.
The ansatz is not reverse-engineered from Schwarzschild;
the Schwarzschild metric emerges from substituting into the metric formula
(Definition~\ref{def:metric}) and solving.

\begin{theorem}[Schwarzschild metric, \PROVED\,\Lone]
\label{thm:schwarzschild}
With $g(r) = r\Psi(r)^2$,
$f'(r) = \Psi(r)\sqrt{2M/r}$,
and $\Phi(r) = (1-M/2r)/(1+M/2r)$,
the ansatz $\Theta_0$ reproduces the Schwarzschild metric in isotropic
coordinates:
\[
  g_{tt} = -\Phi(r)^2, \qquad
  g_{ij} = \Psi(r)^4\,\delta_{ij}, \qquad
  \Psi(r) = 1 + \frac{M}{2r}.
\]
\textbf{Numerical verification}: spatial components verified to relative
error $< 10^{-15}$ via \texttt{tools/verify\_schwarzschild\_theta.py}
(Appendix~\ref{app:numerics}).
\end{theorem}

\begin{tcolorbox}[colback=orange!5!white,colframe=orange!50!black,
  title={Limitation --- temporal component (expected, not an error)}]
The static real-quaternion ansatz has $\partial_t\Theta_0 = 0$, so the metric
formula gives $g_{tt} = 0$ from the spatial construction alone.
The Lorentzian $g_{tt} = -\Phi^2$ is recovered only when the full
complex-time structure $\tau = t + i\psi$ is used:
\[
  \partial_\psi\Theta_0 = i\Phi(r)\Theta_0
  \quad\Longrightarrow\quad
  g_{tt} = -\re[\Tr(\partial_\psi\Theta_0 \cdot \partial_\psi\Theta_0^\dagger)]
          = -\Phi(r)^2.
\]
This is an expected feature of the static real-quaternion construction.
The Lorentzian signature $g_{00} < 0$ is guaranteed by
Theorem~\ref{thm:signature} (AXIOM-B), which operates at the algebraic level.
\end{tcolorbox}

\subsection{ASD Condition and Twistor Space}

\begin{theorem}[ASD condition, \PROVED\,\Lone]
\label{thm:asd}
For $\Theta \in \mathrm{SU}(2)_- \subset \CC\otimes\HH$, $|\Theta| = 1$,
smooth:
\begin{enumerate}
  \item The holonomy of $g_{\mu\nu}[\Theta]$ lies in
        $\mathrm{Sp}(1) \cong \mathrm{SU}(2)_-$, implying the anti-self-dual
        (ASD) Weyl condition $C^+ = 0$.
  \item Combined with the vacuum equation $\nabla^\dagger\nabla\Theta = 0$
        (which gives $R_{\mu\nu} = 0$ via the GR chain), the metric is
        ASD Ricci-flat.
  \item By the Penrose nonlinear graviton theorem~\cite{penrose1976},
        $g_{\mu\nu}[\Theta]$ admits a curved twistor space description.
\end{enumerate}
\emph{Note}: Schwarzschild (Petrov type~D) lies outside the $\mathrm{SU}(2)_-$
sector; it has a non-zero self-dual Weyl tensor.
\end{theorem}

% ==============================================================================
\section{Linearised Gravity and the Regge-Wheeler Equation}
\label{sec:rw}
% ==============================================================================

\begin{theorem}[Linearised GR and Regge-Wheeler, \PROVED\,\Lone]
\label{thm:rw}
Linearising the UBT field equation around flat background
$\Theta = \Theta_0 + \epsilon\,\delta\Theta$
reproduces the linearised Einstein equations.
For odd-parity (axial) perturbations of the Schwarzschild background
decomposed into angular modes $(\ell,m,\omega)$,
the perturbation equation reduces to the
\textbf{Regge-Wheeler equation}~\cite{regge_wheeler1957}:
\[
  \left[\frac{\mathrm{d}^2}{\mathrm{d}r_*^2} + \omega^2
        - V_{\mathrm{RW}}(r)\right]\Psi_{\mathrm{RW}} = 0,
\]
where $r_* = r + 2M\ln|r/2M - 1|$ is the tortoise coordinate and
\[
  V_{\mathrm{RW}}(r)
  = \left(1 - \frac{2M}{r}\right)
    \left[\frac{\ell(\ell+1)}{r^2} - \frac{6M}{r^3}\right]
\]
is the Regge-Wheeler potential for spin-2 gravitational perturbations.
No additional input from UBT beyond the metric chain is required.
\end{theorem}

\begin{tcolorbox}[colback=orange!4!white,colframe=orange!55!black,
  title={Open: Zerilli equation (even-parity graviton) \OPEN\ [\Ltwo]}]
The even-parity (polar) perturbation equation of Schwarzschild
--- the Zerilli equation~\cite{zerilli1970} ---
has not yet been derived from UBT.
This is GAP-Z (\S\ref{sec:gapz}); it does not affect the main GR recovery result.
Closing it requires extending the even-parity $\Theta$ sector analysis
using Chandrasekhar's transformation~\cite{chandrasekhar1983}.
\end{tcolorbox}

% ==============================================================================
\section{Open Problems}
\label{sec:gaps}
% ==============================================================================

The following problems are explicitly bounded.
\textbf{None of them affect the validity of the Main Theorem or
Theorems~\ref{thm:metric}--\ref{thm:rw}.}

\begin{tcolorbox}[colback=red!4!white,colframe=red!55!black,
  title=GAP-10: Off-Shell $\Theta$-Only Closure\ \OPEN\ [\Ltwo],breakable]

\textbf{On-shell result (proved)}: For $\Theta \in \calA_{\mathrm{UBT}}$
satisfying its own Euler--Lagrange equation under a gauge-reduced local rank
condition,
$\delta\hat{S}/\delta\Theta = 0$ is equivalent to the Einstein equations
evaluated on $g = g[\Theta]$.

\smallskip
\textbf{Off-shell gap}: Global non-degeneracy of
$J = \delta g^{\mu\nu}/\delta\Theta$ for \emph{all} field configurations
(not only on-shell $\Theta \in \calA_{\mathrm{UBT}}$) is not proved.

\smallskip
\textbf{Missing lemma}: Show that $\ker J$ consists only of gauge directions
(pure phase or diffeomorphism) for all $\Theta$ in the full off-shell field space.

\smallskip
\textbf{Known obstructions}:
\begin{itemize}
  \item \emph{Rank mismatch}: $\re(\nabla^\dagger\nabla\Theta)$ is rank-0;
        $G_{\mu\nu}$ is rank-2.
  \item \emph{Topology}: global injectivity of $\Theta \to g[\Theta]$ requires
        $H^2(M^4,\mathbb{Z})$ analysis of the $\Theta$-bundle.
  \item \emph{Non-perturbative}: a fixed-point theorem in Sobolev space is needed.
\end{itemize}

\textbf{Scope}: This is a question about off-shell path-integral completeness
and quantum theory.
The classical GR recovery result (all theorems in
\S\S\ref{sec:step1}--\ref{sec:rw}) is unaffected.
\end{tcolorbox}

\subsection{GAP-Z: Zerilli Equation (Even-Parity Graviton)}
\label{sec:gapz}

\begin{tcolorbox}[colback=orange!4!white,colframe=orange!55!black,
  title=GAP-Z: Zerilli Equation\ \OPEN\ [\Ltwo]]
\textbf{Proved}: The Regge-Wheeler equation (odd-parity graviton) is derived
from linearised UBT (Theorem~\ref{thm:rw}).

\smallskip
\textbf{Missing}: The Zerilli equation for even-parity perturbations:
\[
  \left[\frac{\mathrm{d}^2}{\mathrm{d}r_*^2} + \omega^2 - V_\mathrm{Z}(r)\right]
  \Psi_\mathrm{Z} = 0.
\]
Even-parity modes couple scalar and tensor sectors, requiring Chandrasekhar's
two-potential transformation~\cite{chandrasekhar1983}, which has not been
implemented in the UBT even-parity $\Theta$ sector.

\textbf{Priority}: Highest-priority open item for the graviton sector.
Closing strategy: (1) derive the even-parity linearised UBT field equation;
(2) show it reduces to Zerilli via Chandrasekhar's transformation.
\end{tcolorbox}

\subsection{Lower-Priority Gaps}

\begin{center}
\begin{tabular}{lll}
\toprule
\textbf{Gap} & \textbf{Description} & \textbf{Priority} \\
\midrule
GAP-C & FRW/de~Sitter $\Theta$ ansatz not derived & Medium \\
GAP-M & Compact $M^4$ off-shell closure & Low \\
GAP-Q & Path-integral quantisation of UBT & Very long term \\
\bottomrule
\end{tabular}
\end{center}

None of these blocks the on-shell classical result.

% ==============================================================================
\section{Discussion and Comparison to Einstein--Hilbert GR}
\label{sec:discussion}
% ==============================================================================

\subsection{Comparison to Standard GR}

Table~\ref{tab:comparison} gives a systematic comparison between the input
requirements of standard Einstein--Hilbert GR and the UBT derivation.

\begin{table}[ht]
\centering
\caption{Comparison: standard GR inputs versus UBT derivations.}
\label{tab:comparison}
\begin{tabular}{lcc}
\toprule
\textbf{Quantity} & \textbf{Standard GR} & \textbf{UBT (this paper)} \\
\midrule
Spacetime metric $g_{\mu\nu}$ & \emph{Postulated} & \textbf{Derived} (Thm~\ref{thm:metric}) \\
$\det(g) \neq 0$ & \emph{Assumed} & \textbf{Proved} (Thm~\ref{thm:nondeg}) \\
Lorentzian signature $(-,+,+,+)$ & \emph{Assumed} & \textbf{Proved} (Thm~\ref{thm:signature}) \\
$G_{\mu\nu} = 8\pi G\,T_{\mu\nu}$ & \emph{Postulated (EH action)} & \textbf{Derived} (Thm~\ref{thm:einstein}) \\
$T_{\mu\nu}$ symmetric & \emph{Assumed} & \textbf{Proved} \\
$\nabla^\mu T_{\mu\nu} = 0$ & Bianchi + diff.\ inv. & \textbf{Proved} (same) \\
Schwarzschild solution & Solved separately & \textbf{Derived} from $\Theta_0$ (Thm~\ref{thm:schwarzschild}) \\
Regge-Wheeler eq. & Separate derivation & \textbf{Derived} (Thm~\ref{thm:rw}) \\
\bottomrule
\end{tabular}
\end{table}

The Einstein equations emerge in UBT from the same variational principle as
in standard GR (Hilbert action), but the \emph{metric itself} is not an
independent degree of freedom --- it is a bilinear of the fundamental field
$\Theta$.
The Einstein--Hilbert action $S_{\mathrm{EH}} = (16\pi G)^{-1}\int\!\sqrt{-g}\,R\,\mathrm{d}^4x$
appears in the total UBT action as the gravitational term; the $\Theta$-matter
term $S_\Theta$ provides the stress-energy coupling.
The standard GR derivation of $G_{\mu\nu} = 8\pi G\,T_{\mu\nu}$ is
reproduced identically inside UBT.

\subsection{Relation to Existing Algebraic Approaches}

Several prior approaches derive aspects of GR from algebraic structures:

\begin{itemize}
  \item \textbf{Twistor theory}~\cite{penrose1976,penrose_rindler1984}:
        uses the two-spinor calculus of $\mathrm{SL}(2,\CC)$ to reformulate
        GR; the metric is still a separate object.
        UBT's ASD sector (Theorem~\ref{thm:asd}) connects directly to twistor
        theory via the Penrose nonlinear graviton construction.

  \item \textbf{Loop Quantum Gravity}~\cite{thiemann2007}:
        the spin-network state space quantises the metric, but GR is
        recovered as a classical limit, not derived.

  \item \textbf{Connes--Lott noncommutative geometry}~\cite{connes1991}:
        derives a Dirac operator from a spectral triple over a
        21-real-dimensional algebra; a Lorentzian metric is recovered
        from the Dirac operator.
        UBT uses an 8-real-dimensional algebra $\BB$ and derives the
        metric from the bilinear of the fundamental field.

  \item \textbf{Prior biquaternion gravity}~\cite{adler1995,finkelstein1962,deleo1996}:
        these papers work in $\HH$ or $\BB$ but postulate or impose the
        metric.
        The novel step in UBT is the bilinear metric formula
        $g_{\mu\nu} = \re[\Tr(\partial_\mu\Theta \cdot \partial_\nu\Theta^\dagger)]/\calN$,
        which has no analogue in this literature.
\end{itemize}

\subsection{Scope of This Paper}

This paper establishes the \emph{classical GR sector} of UBT.
It makes no claim about the quantum gravity sector (GAP-Q),
cosmological solutions (GAP-C), or gauge unification (T2\_GAUGE track).
The two [L2] open problems (GAP-10 and GAP-Z) are precisely bounded in
\S\ref{sec:gaps} and do not affect the validity of Theorems~\ref{thm:metric}--\ref{thm:rw}.

\subsection{Conclusion}

The chain
\[
  \Theta \;\longrightarrow\;
  g_{\mu\nu} \;\longrightarrow\;
  \Gamma^\lambda_{\mu\nu} \;\longrightarrow\;
  R^\rho{}_{\sigma\mu\nu} \;\longrightarrow\;
  G_{\mu\nu} = 8\pi G\,T_{\mu\nu}
\]
is complete at the [L1] level.
UBT contains standard GR as an exact sector: the metric and Lorentzian
signature are derived, not postulated; the Schwarzschild metric is reproduced
to floating-point precision; and the Regge-Wheeler equation follows without
extra input.

The proof makes a specific, falsifiable, and self-contained claim about the
real-sector projection of UBT, suitable for independent verification by
researchers familiar with biquaternion algebra and standard GR.

% ==============================================================================
\section{Proof Status Summary}
\label{sec:summary}
% ==============================================================================

\begin{tcolorbox}[colback=blue!3!white,colframe=blue!55!black,
  title={T1\_GR Proof Status: Complete (on-shell classical GR)}]
{\small
\begin{tabularx}{\linewidth}{clXl}
\toprule
\textbf{Step} & \textbf{Claim} & \textbf{Status} & \textbf{Source} \\
\midrule
1 & Metric $g_{\mu\nu}$ from $\Theta$ & Proved \Lone & \texttt{step1\_metric} \\
2 & Non-degeneracy $\det(g)\neq 0$ & Proved \Lone & \texttt{step2\_nondeg.} \\
3 & Signature $(-,+,+,+)$ from AXIOM-B & Proved \Lone & \texttt{step3\_sig.} \\
4 & $g \to \Gamma \to R$ chain & Standard GR & \cite{wald1984,mtw1973} \\
5 & Einstein eqs $G_{\mu\nu}=8\pi GT_{\mu\nu}$ & Proved \Lone & \texttt{step3\_einstein} \\
6a & Schwarzschild metric & Proved \Lone & \texttt{verify\_schwarz.py} \\
6b & ASD condition / twistor & Proved \Lone & \texttt{asd\_condition} \\
6c & Regge-Wheeler equation & Proved \Lone & linearised GR chain \\
7a & Zerilli equation (even-parity) & Open \Ltwo & GAP-Z, \S\ref{sec:gapz} \\
7b & Off-shell $\Theta$-only closure & Open \Ltwo & GAP-10 \\
\bottomrule
\end{tabularx}
}
\end{tcolorbox}

% ==============================================================================
% APPENDICES
% ==============================================================================
\appendix

% ==============================================================================
\section{Full Proof of the Signature Theorem}
\label{app:signature}
% ==============================================================================

This appendix gives the full algebraic argument for
Theorem~\ref{thm:signature} (Lorentzian signature from AXIOM-B).

\paragraph{Step A1: Clifford identification.}
The biquaternion algebra satisfies
$\BB = \CC\otimes_\RR\HH \cong \mathrm{Cl}_{1,3}(\RR)$
as real algebras~\cite{porteous1995}.
Under this isomorphism the generators split into one timelike generator
$\gamma^0$ and three spacelike generators $\gamma^i$ satisfying
\[
  (\gamma^0)^2 = -1, \qquad
  (\gamma^i)^2 = +1, \qquad
  \{\gamma^\mu,\gamma^\nu\} = 2\eta^{\mu\nu}.
\]

\paragraph{Step A2: AXIOM-B constraint.}
AXIOM-B states that $\partial_\tau$ lies in the timelike sector of
$\mathrm{Cl}_{1,3}(\RR)$.
Concretely: $\partial_\tau = \partial_t + i\partial_\psi$ satisfies
$\langle\partial_\tau,\partial_\tau\rangle_\eta < 0$.

\paragraph{Step A3: Temporal metric component.}
\[
  g_{00}
  = \frac{\re[\Tr(\partial_0\Theta \cdot \partial_0\Theta^\dagger)]}{\calN}.
\]
With $\partial_0 = \partial_t$ in the timelike $\gamma^0$-sector,
$\partial_0\Theta$ lies in the timelike Clifford sector, so
$\re[\Tr(\partial_0\Theta \cdot \partial_0\Theta^\dagger)] < 0$.
Choosing $\calN = -\re[\Tr(\partial_0\Theta \cdot \partial_0\Theta^\dagger)] > 0$
gives $g_{00} = -1 < 0$.

\paragraph{Step A4: Spatial metric components.}
For $i = 1,2,3$, $\partial_i\Theta$ lies in the spacelike $\gamma^i$-sector,
so $\re[\Tr(\partial_i\Theta \cdot \partial_i\Theta^\dagger)] > 0$,
giving $g_{ii} > 0$.

\paragraph{Step A5: Off-diagonal components.}
Off-diagonal components $g_{0i}$ and $g_{ij}$ for $i \neq j$ arise from
mixed products of time- and space-like Clifford elements.
For the static spatially isotropic sector these vanish by orthogonality
of $\gamma^0$ and $\gamma^i$ under the Clifford inner product.
General spacetimes may have $g_{0i} \neq 0$ (frame dragging), consistent
with the Lorentzian block structure $(-;+,+,+)$.

\paragraph{Conclusion.}
The eigenvalue structure of $g_{\mu\nu}$ has exactly one negative eigenvalue
and three positive eigenvalues, i.e.\ signature $(-,+,+,+)$.
This follows entirely from AXIOM-B, independently of the specific form of $\Theta$.

% ==============================================================================
\section{Numerical Verification of the Schwarzschild Metric}
\label{app:numerics}
% ==============================================================================

Script: \texttt{tools/verify\_schwarzschild\_theta.py}.
Mass $M = 1$ (geometrised units).

\paragraph{ODE consistency condition.}
The ansatz functions satisfy $f'^2 + g'^2 = \Psi^4$ (proved analytically):
\[
  f'^2 + g'^2
  = \Psi^2\,\frac{2M}{r} + \left(1 - \frac{M^2}{4r^2}\right)^2
  = \Psi^4. \qquad \checkmark
\]

\begin{center}
\begin{tabular}{rrrrrr}
\toprule
$r/M$ & $f'$ & $g'$ & $f'^2+g'^2$ & $\Psi^4$ & error \\
\midrule
2.0   & 1.250000 & 0.937500 & 2.441406 & 2.441406 & $0.00$               \\
5.0   & 0.695701 & 0.990000 & 1.464100 & 1.464100 & $1.5\times10^{-16}$  \\
10.0  & 0.469574 & 0.997500 & 1.215506 & 1.215506 & $1.8\times10^{-16}$  \\
100.0 & 0.142128 & 0.999975 & 1.020151 & 1.020151 & $4.3\times10^{-16}$  \\
\bottomrule
\end{tabular}
\end{center}

\paragraph{Spatial metric components.}

\begin{center}
\begin{tabular}{rrrrrrc}
\toprule
$r/M$ & $\Psi^4$ & $g_{xx}$ & $g_{yy}$ & $g_{zz}$ & $\max|g_\mathrm{off}|$ & Status \\
\midrule
2.0   & 2.441406 & 2.441406 & 2.441406 & 2.441406 & $5.6\times10^{-17}$ & OK \\
5.0   & 1.464100 & 1.464100 & 1.464100 & 1.464100 & $1.9\times10^{-16}$ & OK \\
10.0  & 1.215506 & 1.215506 & 1.215506 & 1.215506 & $2.8\times10^{-17}$ & OK \\
50.0  & 1.040604 & 1.040604 & 1.040604 & 1.040604 & $7.3\times10^{-17}$ & OK \\
100.0 & 1.020151 & 1.020151 & 1.020151 & 1.020151 & $2.7\times10^{-16}$ & OK \\
\bottomrule
\end{tabular}
\end{center}

All spatial components agree with $\Psi^4\delta_{ij}$ to floating-point precision.
Off-diagonal components are numerically zero (spherical symmetry).

\paragraph{Honest accounting.}
Spatial components $g_{ij} = \Psi^4\delta_{ij}$: \textbf{verified} (exact + numerical).
Off-diagonal $g_{i0} = 0$: \textbf{verified} (static, spherically symmetric).
Temporal component $g_{tt} = -\Phi^2$: analytically understood via complex-time
structure ($\partial_\psi\Theta_0 = i\Phi\Theta_0$); numerical verification
requires the complex-time UBT solver (planned as future work).

\paragraph{Reproduction.}
\begin{verbatim}
# Prerequisites: pip install numpy
# Run from repository root
python tools/verify_schwarzschild_theta.py
# With optional arguments:
python tools/verify_schwarzschild_theta.py --mass 2.0 --r_values 3,6,12
\end{verbatim}
Exit code 0: all components agree within tolerance $10^{-8}$.
Exit code 1: failure.

% ==============================================================================
% Bibliography
% ==============================================================================
\bibliographystyle{unsrtnat}
\bibliography{UBT_GR_Flagship}

\end{document}
